% P15: signed measures, the three decompositions and density representation.
\chapter{Signed Measures and Radon--Nikodym Densities}
\label{sec:p15-radon-nikodym}

An integral constructs a measure from a density. The converse asks when
one given measure can be represented by integrating against another.
This chapter proves the converse under its precise hypotheses, and
separates the part admitting a density from a singular remainder.
The source is Durrett's Appendix~A.4, including its decomposition
proofs and exercises \cite[pp.~466--470]{DurrettPTEAuthor2019}.
The integral, Stieltjes measures and countable independent coordinates
used below have already been constructed.

\section{Signed measures and positive and negative sets}
\label{sec:p15-signed-measures}

\begin{hbdefinition}{def:p15-signed-measure}{The source's one-sided extended signed-measure convention}
On a measurable space \((\Omega,\mathcal F)\), a signed measure in this
chapter is a map \(\alpha:\mathcal F\to(-\infty,\infty]\) with
\(\alpha(\varnothing)=0\) and countable additivity on disjoint measurable
unions. More precisely, if \(E=\bigcup_i E_i\) is disjoint, then
\[
                         \alpha(E)=\sum_i\alpha(E_i),
\]
where, if \(\alpha(E)\) is finite, the series converges absolutely,
and, if \(\alpha(E)=\infty\), its negative-part sum is finite and its
positive-part sum is infinite. For a scalar \(t\), the parts are
\(t^+=\max(t,0)\) and \(t^-=\max(-t,0)\).
Thus no undefined difference of two infinite values is allowed.
A finite signed measure takes only finite real values.

This follows the source's convention excluding \(-\infty\).
Negating a signed measure that takes \(+\infty\) belongs to the other
one-sided convention, not to the convention just defined.
Nonnegative measures are called measures without the adjective signed.
\index{signed measure}
\end{hbdefinition}

\begin{hbdefinition}{def:p15-signed-positive-sets}{Positive, negative and signed-null sets}
A measurable set \(A\) is positive for \(\alpha\) if
\(\alpha(E)\geq0\) for every measurable \(E\subseteq A\).
It is negative if every such value is at most zero.
It is signed-null if every such value is zero.
The single equality \(\alpha(A)=0\) does not by itself make \(A\)
signed-null, because positive and negative contributions can cancel.
\end{hbdefinition}

\begin{hbtheorem}{thm:p15-signed-measure-examples}{proved=annex:p15.signed-measure-examples}{Differences of measures and integral-defined signed measures}
If \(\nu_1,\nu_2\) are measures with \(\nu_2(\Omega)<\infty\), then
\(\alpha(E)=\nu_1(E)-\nu_2(E)\) is a signed measure in this convention.
If \(f\) is an extended real measurable function with
\(\int f^-\,d\mu<\infty\), then
\[
                            \alpha(E)=\int_E f\,d\mu
\]
is such a signed measure. For this integral-defined example,
\[
 \begin{aligned}
 A\text{ is positive}&\quad\Longleftrightarrow\quad
                           \mu(A\cap\{f<0\})=0,\\
 A\text{ is negative}&\quad\Longleftrightarrow\quad
                           \mu(A\cap\{f>0\})=0.
 \end{aligned}
\]
This supplies Examples A.4.1--2 and Exercise A.4.1, including the
distinction between a zero total value and a signed-null set.
\end{hbtheorem}

\begin{hblemma}{thm:p15-signed-set-closure}{proved=annex:p15.signed-set-closure}{Closure of positive and negative sets}
Every measurable subset of a positive set is positive, and a countable
union of positive sets is positive. The corresponding two assertions
hold for negative sets. Every negative set has finite signed measure.
\end{hblemma}

\begin{hblemma}{thm:p15-negative-subset}{proved=annex:p15.negative-subset}{A negative-valued event contains a negative set}
If \(E\in\mathcal F\) and \(\alpha(E)<0\), there is a negative measurable
set \(F\subseteq E\) with \(\alpha(F)\leq\alpha(E)<0\).
The annex gives the source's successive positive-set removal and handles
both finite termination and its infinite remainder.
\end{hblemma}

\section{Hahn and Jordan decomposition}
\label{sec:p15-hahn-jordan}

\begin{hbtheorem}{thm:p15-hahn-decomposition}{proved=annex:p15.hahn-decomposition}{Hahn decomposition and its exact uniqueness}
Every signed measure in the stated convention has a partition
\(\Omega=A\mathbin{\dot\cup}B\) with \(A\) positive and \(B\) negative.
The negative set can be chosen to attain
\[
                    \alpha(B)=\inf\{\alpha(N):N\text{ negative}\}>-\infty.
\]
If \((A_1,B_1)\) and \((A_2,B_2)\) are two such partitions, then
\(A_1\cap B_2\) and \(A_2\cap B_1\) are signed-null.
The partition itself need not be unique.
\end{hbtheorem}

\begin{hbdefinition}{def:p15-singularity}{Mutual singularity}
Measures \(\mu,\nu\) on the same measurable space are mutually singular,
written \(\mu\perp\nu\), if some \(N\in\mathcal F\) satisfies
\(\mu(N)=0\) and \(\nu(N^c)=0\).
The condition is symmetric by using the complementary set.
Thus \(\nu\) is concentrated on a measurable set that is null for \(\mu\).
\index{mutually singular measures}
\end{hbdefinition}

\begin{hbtheorem}{thm:p15-jordan-decomposition}{proved=annex:p15.jordan-decomposition}{Jordan components and total variation}
There is a unique pair of mutually singular measures
\(\alpha^+,\alpha^-\) with
\[
                 \alpha=\alpha^+-\alpha^-,
                 \qquad \alpha^-(\Omega)<\infty.
\]
For any Hahn partition \((A,B)\), they are
\[
           \alpha^+(E)=\alpha(E\cap A),\qquad
           \alpha^-(E)=-\alpha(E\cap B).
\]
Define the total variation measure by
\(\lvert\alpha\rvert=\alpha^++\alpha^-\).
The bars here name a measure and are not set cardinality.
For every measurable \(E\),
\[
 \begin{aligned}
 \alpha^+(E)&=\sup_{F\in\mathcal F,\ F\subseteq E}\alpha(F),\\
 \alpha^-(E)&=-\inf_{F\in\mathcal F,\ F\subseteq E}\alpha(F),\\
 \lvert\alpha\rvert(E)
   &=\sup_{E=\mathbin{\dot\bigcup}_{j=1}^m E_j}
                         \sum_{j=1}^m|\alpha(E_j)|.
 \end{aligned}
\]
The last supremum ranges over finite measurable partitions and permits
infinite values. A set is signed-null exactly when its total variation
is zero. The signed measure is finite exactly when
\(\lvert\alpha\rvert(\Omega)<\infty\).
\index{Jordan decomposition}\index{total variation!signed measure}
\end{hbtheorem}

\begin{hbtheorem}{thm:p15-signed-cdf-variation}{proved=annex:p15.signed-cdf-variation}{Finite signed distribution functions and interval variation}
For a finite signed Borel measure on \(\mathbb R\), put
\[
 H(x)=\alpha((-\infty,x]),\quad
 F(x)=\alpha^+((-\infty,x]),\quad
 G(x)=\alpha^-((-\infty,x]).
\]
Then \(H=F-G\), where \(F,G\) are bounded, nonnegative, nondecreasing,
right continuous and tend to zero at \(-\infty\).
Conversely any two functions with those properties determine a unique
finite signed Borel measure with cumulative function \(F-G\).
Their individual Stieltjes measures need not already be the mutually
singular Jordan components.

For finite real \(a<b\),
\[
 \lvert\alpha\rvert((a,b])
  =\sup_{a=t_0<\cdots<t_m=b}
       \sum_{j=1}^m|H(t_j)-H(t_{j-1})|.
\]
The initial jump at \(a\) is excluded by the half-open measure interval.
The proof supplies the finite-measure interval approximation needed
for the reverse inequality.
\end{hbtheorem}

\section{The absolutely continuous and singular parts}
\label{sec:p15-lebesgue-decomposition}

\begin{hbdefinition}{def:p15-absolute-continuity}{Absolute continuity of a measure}
Write \(\nu\ll\mu\) when
\[
                     \mu(E)=0\quad\Longrightarrow\quad\nu(E)=0
                              \qquad(E\in\mathcal F).
\]
This is a relation between measures, not a pointwise continuity
condition on a density. For a signed measure, \(\alpha\ll\mu\) means
that the same implication holds with \(\alpha(E)=0\).
All subsets considered remain measurable in the original sigma-algebra.
\index{absolute continuity!measures}
\end{hbdefinition}

\begin{hbtheorem}{thm:p15-lebesgue-decomposition}{proved=annex:p15.lebesgue-decomposition}{Lebesgue decomposition for sigma-finite measures}
Let \(\mu,\nu\) be sigma-finite measures on \((\Omega,\mathcal F)\).
There exist unique measures \(\nu_{\mathrm r},\nu_{\mathrm s}\) with
\[
 \nu=\nu_{\mathrm r}+\nu_{\mathrm s},\qquad
 \nu_{\mathrm r}\ll\mu,\qquad \nu_{\mathrm s}\perp\mu.
\]
The regular part has a nonnegative measurable density:
\[
                         \nu_{\mathrm r}(E)=\int_E g\,d\mu.
\]
It has a version \(g\) finite \(\mu\)-almost everywhere.
The proof first maximizes the integral of a dominated density on a
finite-measure space, proves singularity of the remainder by Hahn
decomposition, and then assembles a common finite partition for the two
sigma-finite measures. It proves uniqueness as well as existence.
\end{hbtheorem}

\begin{hbtheorem}{thm:p15-radon-nikodym}{proved=annex:p15.radon-nikodym}{Radon--Nikodym representation and uniqueness}
If \(\mu,\nu\) are sigma-finite and \(\nu\ll\mu\), there is a
nonnegative measurable \(g\), finite \(\mu\)-almost everywhere, with
\[
                              \nu(E)=\int_E g\,d\mu
                                      \qquad(E\in\mathcal F).
\]
Any two such densities agree \(\mu\)-almost everywhere.
A representative is denoted \(d\nu/d\mu\); its values on a measurable
\(\mu\)-null set may be assigned arbitrarily while preserving measurability.
Conversely, any measure given by such an integral is absolutely continuous
with respect to \(\mu\).
\index{Radon--Nikodym derivative}
\end{hbtheorem}

\begin{hbtheorem}{thm:p15-density-calculus}{proved=annex:p15.density-calculus}{Change of measure, sums, chains and reciprocals}
All measures in the derivative statements below are sigma-finite.
\begin{enumerate}
\item If \(\nu\ll\mu\), \(g=d\nu/d\mu\), and \(f\geq0\) is measurable, then
      \[
                        \int f\,d\nu=\int fg\,d\mu,
      \]
      with the previously defined nonnegative convention \(0\cdot\infty=0\).
      The same equality holds for signed measurable \(f\) whenever its signed
      integral is defined; corresponding positive and negative integrals agree.
\item If \(\nu_1,\nu_2\ll\mu\), then
      \[
               \frac{d(\nu_1+\nu_2)}{d\mu}
                  =\frac{d\nu_1}{d\mu}+\frac{d\nu_2}{d\mu}
                          \quad\mu\text{-a.e.}
      \]
      Nonnegative finite scalar multiples obey the corresponding rule.
\item If \(\pi\ll\nu\ll\mu\), then
      \[
                  \frac{d\pi}{d\mu}
                    =\frac{d\pi}{d\nu}\frac{d\nu}{d\mu}
                         \quad\mu\text{-a.e.}
      \]
\item If \(\nu\ll\mu\ll\nu\), then \(0<d\nu/d\mu<\infty\)
      \(\mu\)-a.e. and
      \[
                 \frac{d\mu}{d\nu}=\left(\frac{d\nu}{d\mu}\right)^{-1}
                         \quad\nu\text{-a.e.}
      \]
      The reciprocal may be defined as zero off the positive finite set.
\item If \(\mu_1\ll\mu_2\) and \(\mu_2\perp\nu\), then \(\mu_1\perp\nu\).
      This last assertion does not require sigma-finiteness.
\end{enumerate}
The proof specifies all null-set qualifications in the chain and reciprocal
rules and supplies Exercises A.4.5--9.
\end{hbtheorem}

\begin{hbcorollary}{thm:p15-finite-signed-density}{proved=annex:p15.finite-signed-density}{Densities of finite signed measures}
If \(\alpha\) is a finite signed measure, \(\mu\) is sigma-finite and
\(\alpha\ll\mu\), then some integrable real measurable \(f\) satisfies
\[
                            \alpha(E)=\int_E f\,d\mu.
\]
It is unique \(\mu\)-a.e., and
\[
        \frac{d\alpha^+}{d\mu}=f^+,\qquad
        \frac{d\alpha^-}{d\mu}=f^-,\qquad
        \frac{d\lvert\alpha\rvert}{d\mu}=|f|\quad\mu\text{-a.e.}
\]
In particular \(\lvert\alpha\rvert(\Omega)=\int|f|\,d\mu\).
\end{hbcorollary}

\begin{hbtheorem}{thm:p15-finite-uniform-absolute-continuity}{proved=annex:p15.finite-uniform-absolute-continuity}{The epsilon-delta form for a finite target measure}
If \(\nu\) is finite and \(\mu\) is any measure, then \(\nu\ll\mu\)
exactly when for every \(\varepsilon>0\) there is \(\delta>0\) such that
\[
                    \mu(E)<\delta\quad\Longrightarrow\quad
                    \nu(E)<\varepsilon\qquad(E\in\mathcal F).
\]
This assertion is about the mass of all small measurable sets. The finite
target hypothesis is used in the proof's continuity-from-above step.
\end{hbtheorem}

\section{Singular continuous laws and worked decompositions}
\label{sec:p15-density-examples}

\begin{hbdefinition}{def:p15-cantor-law-construction}{The Cantor set and its fair-digit law}
For a binary word \(b=(b_1,\ldots,b_n)\), put
\[
 c_b=\sum_{j=1}^n\frac{2b_j}{3^j},\qquad
 I_b=[c_b,c_b+3^{-n}],\qquad
 C_n=\bigcup_{b\in\{0,1\}^n}I_b,\qquad C=\bigcap_{n\geq1}C_n.
\]
These are the usual middle-thirds stages and Cantor set.
On the constructed probability space with independent fair bits \(B_j\), set
\[
                   X=\sum_{j\geq1}\frac{2B_j}{3^j},\qquad
                   \eta(E)=\mathbb P(X\in E).
\]
The series is bounded by one. The next theorem proves measurability
and the claimed distribution, rather than presuming that a continuous
cumulative function must have a density.
\end{hbdefinition}

\begin{hbtheorem}{thm:p15-cantor-singular-law}{proved=annex:p15.cantor-singular-law}{A continuous distribution with no Lebesgue density}
The random variable \(X\) is measurable, \(\eta(C)=1\),
\(\lambda(C)=0\), and \(\eta(I_b)=2^{-n}\) for every length-\(n\) word.
Every singleton has \(\eta\)-mass zero. Its cumulative function is
continuous, is zero for \(x\leq0\), one for \(x\geq1\), and takes values
\(1/2\) on \([1/3,2/3]\), \(1/4\) on \([1/9,2/9]\), and \(3/4\)
on \([7/9,8/9]\). More generally each removed middle-third interval
has the corresponding dyadic plateau value proved in the annex.
Thus this is the source's uniform Cantor distribution
\cite[example~1.2.7, pp.~13--14; exercise~A.4.2, p.~468]{DurrettPTEAuthor2019}.
It is singular with respect to Lebesgue measure and has no Lebesgue density.
It is not concentrated on any countable set.
\end{hbtheorem}

\begin{hbtheorem}{thm:p15-radon-nikodym-examples}{proved=annex:p15.radon-nikodym-examples}{Explicit derivatives, singular components and failed hypotheses}
The following statements hold.
\begin{enumerate}
\item On \([0,1]\), \(\nu(E)=\int_E2x\,d\lambda+\delta_0(E)\)
      has regular part with density \(2x\) and singular part \(\delta_0\).
      For \(\mu=\lambda+\delta_0\), it is entirely absolutely continuous,
      with \(d\nu/d\mu=2x\) for \(x>0\) and value 1 at zero.
\item For \(\alpha(E)=\int_E(2x-1)\,d\lambda\) on \([0,1]\),
      \(\alpha(\Omega)=0\), but
      \(\alpha^+(\Omega)=\alpha^-(\Omega)=1/4\) and
      \(\lvert\alpha\rvert(\Omega)=1/2\).
\item The source's counting measure on the Borel real line is not
      sigma-finite. Lebesgue measure is absolutely continuous with respect
      to it but has no density with respect to it.
\item On \((0,1)\), the measure with density \(1/x\) relative to
      Lebesgue measure is sigma-finite and absolutely continuous, but
      fails the preceding uniform epsilon-delta assertion.
\end{enumerate}
The base measure is part of every density claim: changing it can change
which part is singular and the value assigned at a formerly null point.
\end{hbtheorem}

\begin{proofboundary}{def:p15-radon-nikodym-next}{Conditional expectations and conditional laws}
The density representation supplies the existence tool for
\cref{sec:p15-conditional-expectation}, which proves the scalar
conditional calculus from its measurability and all-event identity.
\Cref{sec:p15-regular-conditional-laws} constructs regular conditional
kernels and kernels on observation spaces, with their Borel target
hypotheses and common full-measure domains.
The distinct density-martingale construction of the Lebesgue decomposition
is \cref{thm:p15-density-lebesgue-decomposition}; its application to
infinite independent products is \cref{thm:p15-product-affinity-criterion}.
For a countably generated sigma-field, the independent finite-partition
construction in \cref{thm:p15-countably-generated-radon-nikodym}
proves density existence without assuming the general conditional
expectation existence theorem.
\end{proofboundary}
